\documentclass[11pt]{article}
\usepackage[utf8]{inputenc}
\usepackage[T1]{fontenc}
\usepackage{amsmath, amssymb, amsthm, mathtools}
\usepackage{geometry}
\geometry{margin=1in}
\usepackage{hyperref}

\newtheorem{theorem}{Theorem}
\newtheorem{lemma}[theorem]{Lemma}
\newtheorem{proposition}[theorem]{Proposition}
\theoremstyle{definition}
\newtheorem{definition}[theorem]{Definition}
\newtheorem{example}[theorem]{Example}
\theoremstyle{remark}
\newtheorem{remark}[theorem]{Remark}

\newcommand{\R}{\mathbb{R}}
\newcommand{\N}{\mathbb{N}}
\DeclareMathOperator*{\esssup}{ess\,sup}

\title{Existence and Uniqueness for ODEs:\\ The Picard-Lindel\"of Theorem}
\author{Math Foundations --- Node 30}
\date{}

\begin{document}
\maketitle

\subsection*{First, an example}

Before the general theorem, look at the simplest nontrivial IVP:
$y' = y$ with $y(0) = 1$. Here $f(t, y) = y$, so $|f(t, y_1) - f(t, y_2)|
= |y_1 - y_2|$ and $f$ is Lipschitz in $y$ with constant $L = 1$. The
Picard iterates starting from $y_0(t) \equiv 1$ are
\[
  y_1(t) = 1 + \int_0^t 1 \, ds = 1 + t, \qquad
  y_2(t) = 1 + \int_0^t (1 + s) \, ds = 1 + t + \tfrac{t^2}{2},
\]
\emph{Read aloud:} ``y-one of t equals one plus the integral from zero to t
of one d-s, which is one plus t; y-two of t equals one plus the integral
from zero to t of one plus s d-s, which is one plus t plus t-squared over
two.''
and inductively $y_k(t) = \sum_{j=0}^{k} t^j / j!$, the partial sums of
the Taylor series of $e^t$. They converge uniformly on every compact
interval to the unique solution $y(t) = e^t$.

Contrast this with $y' = y^{2/3}$, $y(0) = 0$. Now $\partial_y f = \tfrac23
y^{-1/3}$ blows up at $0$, so $f$ is \emph{not} Lipschitz near $y = 0$,
and the IVP admits both $y \equiv 0$ and $y(t) = (t/3)^3$ as solutions.
Lipschitz continuity is exactly what rules out this branching.

\section{Setting}

Let $D \subset \R \times \R^n$ be an open set, $(t_0, y_0) \in D$, and
$f : D \to \R^n$ a continuous function. We study the initial value problem
(IVP)
\begin{equation}
  \label{eq:ivp}
  y'(t) = f(t, y(t)), \qquad y(t_0) = y_0.
\end{equation}
\emph{Read aloud:} ``y-prime of t equals f of t and y of t, with y of
t-naught equal to y-naught.''
A \emph{solution} is a differentiable curve $y : I \to \R^n$ on an open
interval $I \ni t_0$ satisfying \eqref{eq:ivp} pointwise.

\section{Lipschitz continuity and Picard iteration}

\begin{definition}[Lipschitz in $y$]
  A function $f : D \to \R^n$ is \emph{Lipschitz continuous in $y$} on
  $D$ if there is $L \ge 0$ such that
  \[
    \| f(t, y_1) - f(t, y_2) \| \le L \, \| y_1 - y_2 \|
    \qquad \forall (t, y_1), (t, y_2) \in D.
  \]
  \emph{Read aloud:} ``the norm of f of t y-one minus f of t y-two is at
  most L times the norm of y-one minus y-two, for all pairs in D.''
  The smallest such $L$ is the \emph{Lipschitz constant} of $f$.
\end{definition}

\begin{remark}
  If $f$ is $C^1$ on $D$, then $\partial f / \partial y$ being bounded by
  $L$ on $D$ (in operator norm) implies $f$ is Lipschitz in $y$ with
  constant $L$. This is the most common way Lipschitz hypotheses arise
  in practice.
\end{remark}

\begin{definition}[Picard iteration]
  Given the IVP \eqref{eq:ivp}, define the sequence of functions
  $\{y_k\}_{k \ge 0}$ by
  \[
    y_0(t) \equiv y_0, \qquad
    y_{k+1}(t) \,=\, y_0 + \int_{t_0}^t f\bigl(s, y_k(s)\bigr)\, ds.
  \]
  \emph{Read aloud:} ``y-naught of t is identically y-naught; y-sub-k-plus-one
  of t equals y-naught plus the integral from t-naught to t of f of s and
  y-k of s, d-s.''
  These are the \emph{Picard iterates} of the IVP.
\end{definition}

A continuous function $y$ solves \eqref{eq:ivp} if and only if it is a
fixed point of the integral operator
$T[y](t) = y_0 + \int_{t_0}^t f(s, y(s))\, ds$. So existence and uniqueness
become a fixed-point problem for $T$.

\section{The Picard-Lindel\"of theorem}

\begin{theorem}[Picard-Lindel\"of]
  \label{thm:pl}
  Let $a, b > 0$ and let
  \[
    R = \{ (t, y) \in \R \times \R^n : |t - t_0| \le a, \; \|y - y_0\| \le b \}.
  \]
  \emph{Read aloud:} ``R is the set of pairs t comma y in R-cross-R-n such
  that the absolute value of t minus t-naught is at most a, and the norm
  of y minus y-naught is at most b.''
  Suppose $f$ is continuous on $R$ with $\|f(t, y)\| \le M$ for all
  $(t,y) \in R$, and Lipschitz in $y$ on $R$ with constant $L$. Set
  \[
    \alpha = \min\!\left(a, \, \tfrac{b}{M}\right).
  \]
  \emph{Read aloud:} ``alpha is the minimum of a and b-over-M.''
  Then the IVP \eqref{eq:ivp} has a unique solution
  $y : [t_0 - \alpha, \, t_0 + \alpha] \to \R^n$, and the Picard iterates
  $y_k$ converge to it uniformly on that interval.
\end{theorem}

\begin{proof}[Proof sketch]
  Set $I = [t_0 - \alpha, t_0 + \alpha]$ and let
  \[
    X = \{ y \in C(I, \R^n) : \|y(t) - y_0\| \le b \text{ for all } t \in I \}.
  \]
  \emph{Read aloud:} ``X is the set of continuous functions y from I into
  R-n such that the norm of y of t minus y-naught is at most b for all t
  in I.''
  Equip $C(I, \R^n)$ with the sup norm $\|y\|_\infty = \sup_{t \in I}\|y(t)\|$.
  Then $X$ is a closed subset of a Banach space, hence complete. Define
  $T : X \to C(I, \R^n)$ by
  \[
    T[y](t) = y_0 + \int_{t_0}^t f\bigl(s, y(s)\bigr)\, ds.
  \]
  \emph{Read aloud:} ``T applied to y, evaluated at t, equals y-naught plus
  the integral from t-naught to t of f of s and y of s, d-s.''

  \textbf{Step 1 ($T$ maps $X$ into $X$).} For $y \in X$ and $t \in I$,
  \[
    \|T[y](t) - y_0\| \le \int_{t_0}^t \|f(s, y(s))\|\, ds
    \le M |t - t_0| \le M\alpha \le b.
  \]
  \emph{Read aloud:} ``the norm of T-of-y at t minus y-naught is at most
  the integral of the norm of f, which is at most M times t minus t-naught,
  which is at most M-alpha, which is at most b.''
  Hence $T[y] \in X$.

  \textbf{Step 2 (Contraction in a weighted norm).} On $C(I, \R^n)$ define
  the equivalent norm
  $\|y\|_w = \sup_{t \in I} e^{-2L|t - t_0|}\, \|y(t)\|.$
  Then for $y_1, y_2 \in X$ and $t \ge t_0$ in $I$,
  \begin{align*}
    \|T[y_1](t) - T[y_2](t)\|
      &\le \int_{t_0}^t \|f(s, y_1(s)) - f(s, y_2(s))\| \, ds \\
      &\le L \int_{t_0}^t \|y_1(s) - y_2(s)\| \, ds \\
      &\le L \int_{t_0}^t e^{2L(s-t_0)}\, ds \cdot \|y_1 - y_2\|_w \\
      &= \tfrac{1}{2}\, \bigl(e^{2L(t-t_0)} - 1\bigr)\, \|y_1 - y_2\|_w.
  \end{align*}
  \emph{Read aloud:} ``the norm of T-of-y-one minus T-of-y-two at t is
  bounded by the integral of the difference of f's, then by L times the
  integral of the difference of y's, then by L times the integral of
  e-to-the-two-L-s-minus-t-naught d-s times the weighted norm, equaling
  one-half times e-to-the-two-L-t-minus-t-naught minus one, times the
  weighted norm.''
  Multiplying by $e^{-2L(t-t_0)}$ and taking sup over $t$ (the $t < t_0$
  case is symmetric) yields
  \[
    \|T[y_1] - T[y_2]\|_w \le \tfrac{1}{2}\, \|y_1 - y_2\|_w.
  \]
  \emph{Read aloud:} ``the weighted norm of T-of-y-one minus T-of-y-two is
  at most one-half times the weighted norm of y-one minus y-two.''
  Thus $T$ is a strict contraction on $(X, \|\cdot\|_w)$ with contraction
  factor $1/2$.

  \textbf{Step 3 (Banach fixed point).} By the Banach contraction mapping
  theorem, $T$ has a unique fixed point $y^\ast \in X$. By construction
  $y^\ast$ is continuous and satisfies the integral equation; since $f$
  is continuous, $y^\ast$ is in fact $C^1$ and solves
  \eqref{eq:ivp} pointwise on $I$. Uniqueness in $X$ follows from the
  contraction property; any solution must lie in $X$ (by Step~1 applied
  to it), so uniqueness is global on $I$.

  Finally, the Picard iterates are exactly the sequence $T^k[y_0]$ produced
  by the contraction mapping theorem, and so converge uniformly on $I$
  to $y^\ast$.
\end{proof}

\section{What goes wrong without Lipschitz}

\begin{example}[Non-uniqueness]
  Consider $y' = y^{2/3}$, $y(0) = 0$. Here $f(t, y) = y^{2/3}$ fails to
  be Lipschitz at $y = 0$ since $\partial_y f = \tfrac{2}{3} y^{-1/3}$
  blows up. The IVP admits the two distinct solutions
  $y \equiv 0$ and $y(t) = (t/3)^3$, as well as a one-parameter family
  $y_\tau$ that equals $0$ for $t \le \tau$ and $((t - \tau)/3)^3$ for
  $t > \tau$.
\end{example}

\section{Connection to flow-based generative models}

Modern generative models (flow matching, ELF, neural ODEs) define
\[
  \dot x(t) = v_\theta(t, x(t))
\]
\emph{Read aloud:} ``x-dot of t equals v-sub-theta of t and x of t.''
where $v_\theta$ is a neural network. On any compact set $K$, neural
networks built from smooth activations are $C^1$ and hence Lipschitz in
$x$ with a finite constant $L_K$. Theorem~\ref{thm:pl} therefore
guarantees a unique sample trajectory for each starting point $x_0$, and
hence a well-defined flow $\phi_t$ on $K$. Smooth dependence on initial
conditions further makes $\phi_t$ a diffeomorphism --- a property used
implicitly throughout the literature.

\end{document}
